% !TEX root = ../../tree_notes.tex
% =====================================================================
%  16.4  Lifting an Aggregate along a Jump          sec:c11-aggregates
%  (This section absorbs the material originally planned for old Ch.12
%   "Binary Lifting" whose folder was an unwritten 3-topic outline.)
%  Exports: thm:c11-agg  lst:c11-pathagg  ex:c11-aggrun
%           prob:c11-1304e  lst:c11-1304e-soln
% =====================================================================
\section{Lifting an Aggregate along a Jump}\label{sec:c11-aggregates}

\begin{previouslybox}[Three questions from 16.1--16.3]
\begin{enumerate}[nosep, leftmargin=1.8em]
  \item From \secref{c11-lifting}: the \(\code{up}[v][j]\) table lets you
        jump \(2^j\) ancestors in one step.  \emph{(What if each edge
        carried a \emph{value} and you needed the \(\max\) along the
        jump?  Can the same table structure carry a payload?)}
  \item From Ch.~13 (\(\thmref{c10-subtree-template}\)): the subtree fold
        combines child contributions.  \emph{(Lifting an aggregate along
        a path is the same idea on \emph{ancestor chains}: associative
        combine, partial results, merge at the end.)}
  \item From \(\exref{c11-census}\): \(\lca(8,10) = 1\) splits the path
        \(8\to\cdots\to10\) at the hinge.  \emph{(If you want
        \(\max\) along that path, you can accumulate on each side
        separately, then merge at the hinge.)}
\end{enumerate}
\end{previouslybox}

\begin{topicintro}
The lifting table \(\code{up}[v][j]\) answers ``where am I after
\(2^j\) steps up?''.  A \emph{second} table \(\code{agg}[v][j]\) answers
``what is the aggregate (max, min, sum, xor, count, \(\gcd\), \(\dots\))
along those \(2^j\) steps?''  The machinery is identical -- same
recurrence, same bit-decomposition climb -- with an associative payload
riding on every jump.  Path queries then split at the hinge (\(\lca\)):
accumulate each arm with the \(\code{agg}\) table, merge the two partial
results, and fold in the LCA's own value.  Three distinct patterns emerge
from this one idea:
\begin{itemize}[nosep]
  \item \textbf{Static path aggregate} (max/min/sum/xor along a route);
  \item \textbf{First-ancestor-past-a-threshold} (bit-descent for the
        deepest ancestor whose prefix-sum or value still passes a
        predicate);
  \item \textbf{\(k\)-th node \emph{on} a route} (distance formula
        + \(k\)-th ancestor, the key to CF~1304E).
\end{itemize}

\textbf{After this section you can}
\begin{itemize}[nosep]
  \item define \(\code{agg}[v][j]\) and prove the path-aggregate
        theorem (\thmref{c11-agg}); \hfill\lv{4}
  \item compute max/sum/xor along a path on~\(T\) by hand; \hfill\lv{4}
  \item descend to the first ancestor past a threshold; \hfill\lv{4}--\lv{5}
  \item locate the \(k\)-th node on a route \(u\to v\); \hfill\lv{4}
  \item solve CF~1304E (walk on a tree + one extra edge, parity check).
        \hfill\lv{4}--\lv{5}
\end{itemize}
\end{topicintro}

% ---------------------------------------------------------------
\subsection{The Second Table}
% ---------------------------------------------------------------

\begin{theorem}{Aggregated lifting}{c11-agg}
Let \(\oplus\) be an \emph{associative} binary operation on values
attached to vertices (or edges).  Define
\[
  \code{agg}[v][0] = \mathit{val}(v), \qquad
  \code{agg}[v][j] = \code{agg}[v][j\!-\!1] \oplus
                      \code{agg}\!\bigl[\,\code{up}[v][j\!-\!1]\,\bigr][j\!-\!1].
\]
\(\code{agg}[v][j]\) is the aggregate over the \(2^j\) vertices starting
at~\(v\) going upward, \textbf{excluding} the \emph{landing} vertex
\(\code{up}[v][j]\).  To query the aggregate along the path from~\(u\)
to~\(v\):
\begin{enumerate}[nosep]
  \item Compute \(\ell = \lca(u,v)\).
  \item Accumulate \(u \to \ell\): decompose \(\depth(u)-\depth(\ell)\)
        jumps and \(\oplus\)-combine the partial \(\code{agg}\) values.
  \item Accumulate \(v \to \ell\) symmetrically.
  \item Merge: combine the two partial results and fold in
        \(\mathit{val}(\ell)\) (the hinge vertex, never captured by
        either arm because it is the landing vertex).
\end{enumerate}
Cost \(O(\log n)\) per query.
\end{theorem}

\begin{proofbox}
The recurrence mirrors \(\code{up}\)'s composition: a \(2^j\)-jump from
\(v\) is two consecutive \(2^{j-1}\)-jumps, so the aggregate over the
full \(2^j\) is the combine of the lower half (starting at~\(v\)) and the
upper half (starting at the landing of the lower half).  The "exclude the
landing vertex" convention makes the halves non-overlapping and adjacent.
The query decomposes the path at the hinge; each arm is a \(k\)-ancestor
climb where \(k\) is split into binary bits, and the \(\code{agg}\)
values for those bits tile the arm exactly.  The hinge itself is the only
vertex not covered by either arm's jumps (each arm's last jump lands
\emph{at} the hinge and excludes it), so we fold it in last.
\end{proofbox}

\begin{notebox}[Why \(\oplus\) must be associative -- and when idempotent matters]
The two-accumulator approach (left arm, right arm) \emph{never} overlaps
and \emph{never} misses, so \(\oplus\) need only be associative -- not
commutative or idempotent.  This is \textbf{different} from the
sparse-table overlapping-block trick (\(\S\)\ref{sec:c11-euler}), where
the block overlap requires \(\oplus\) to be \emph{idempotent}
(\(x\oplus x = x\)).  Sums, xors, and counts are fine here but not there.
\end{notebox}

% ---------------------------------------------------------------
\subsection{Worked Example on T}
% ---------------------------------------------------------------

Assign vertex values (permutation of~1..10): \(w = (6,3,9,1,8,10,7,2,5,4)\)
for nodes~1..10 (matching the summary).  \(\oplus = \max\).

\(\code{agg}[v][0] = w_v\) (the vertex's own value):
\((6,3,9,1,8,10,7,2,5,4)\).

\(\code{agg}[v][1] = w_v \oplus w_{\code{par}[v]}\) (max of \(v\) and parent,
excluding the landing vertex, which is the grandparent):
\[
\begin{array}{c|c|c}
v & \code{up}[v][0] = \mathit{par} & \code{agg}[v][1] = \max(w_v, w_{\mathit{par}}) \\ \hline
1 & 0 & 6 \\
2 & 1 & \max(3,6)=6 \\
3 & 1 & \max(9,6)=9 \\
4 & 2 & \max(1,3)=3 \\
5 & 2 & \max(8,3)=8 \\
6 & 3 & \max(10,9)=10 \\
7 & 3 & \max(7,9)=9 \\
8 & 5 & \max(2,8)=8 \\
9 & 7 & \max(5,7)=7 \\
10 & 9 & \max(4,5)=5
\end{array}
\]

\(\code{agg}[v][2] = \code{agg}[v][1] \oplus \code{agg}[\code{up}[v][1]][1]\)
(max over 4 vertices starting at \(v\), excluding the landing):

For \(v = 8\): \(\code{agg}[8][1] = 8\); \(\code{up}[8][1] = 2\);
\(\code{agg}[2][1] = 6\); so \(\code{agg}[8][2] = \max(8,6) = 8\).
For \(v = 10\): \(\code{agg}[10][1] = 5\); \(\code{up}[10][1] = 7\);
\(\code{agg}[7][1] = 9\); so \(\code{agg}[10][2] = \max(5,9) = 9\).

\begin{example}{Max-on-path query on T}{c11-aggrun}
Compute the \(\max\) along the path \(8 \to 10\).
Hinge: \(\lca(8,10) = 1\).

Arm from \(8\): \(\depth(8)-\depth(1) = 3 - 0 = 3\).
Decompose \(3 = 2^1 + 2^0\):
\begin{itemize}[nosep]
  \item \(j\!=\!1\): \(\code{agg}[8][1] = 8\); land at \(\code{up}[8][1]=2\).
  \item \(j\!=\!0\): \(\code{agg}[2][0] = 3\); land at \(\code{up}[2][0]=1\).
\end{itemize}
Left partial: \(\max(8, 3) = 8\).

Arm from \(10\): \(\depth(10)-\depth(1) = 4 - 0 = 4 = 2^2\).
\begin{itemize}[nosep]
  \item \(j\!=\!2\): \(\code{agg}[10][2] = 9\); land at \(\code{up}[10][2]=1\).
\end{itemize}
Right partial: \(9\).

Merge with hinge: \(w_1 = 6\).
Final: \(\max(8, 9, 6) = \boxed{9}\) (at node~3, on the right arm).

Path vertices: \(8, 5, 2, 1, 3, 7, 9, 10\) with values
\((2,8,3,6,9,7,5,4)\); \(\max = 9\).  \checkmark
\end{example}

% ---------------------------------------------------------------
\subsection{Implementation Skeleton}
% ---------------------------------------------------------------

\begin{lstlisting}[caption={General path-aggregate query (max/sum/xor)}, label={lst:c11-pathagg}]
// val[v] = value at vertex v;  ID = identity of op
//  MAX: ID=0 (-inf for negative), op=max, agg[i][j]=max(agg[i][j-1], agg[up[i][j-1]][j-1])
//  SUM: ID=0, op=+
//  XOR: ID=0, op=^

int path_query(int u, int v, int l) {  // l = lca(u,v)
    int res = val[l];   // fold hinge value
    int cur = u;
    int k = depth_[u] - depth_[l];
    for (int j = 0; j < LOG; j++)
        if (k & (1 << j)) { res = op(res, agg[cur][j]); cur = up[cur][j]; }
    cur = v;
    k = depth_[v] - depth_[l];
    for (int j = 0; j < LOG; j++)
        if (k & (1 << j)) { res = op(res, agg[cur][j]); cur = up[cur][j]; }
    return res;
}
\end{lstlisting}

\begin{complexitybox}[Aggregated lifting]
Build: \(O(n\log n)\), same DFS as \(\code{up}\) with one extra array.
Query: \(O(\log n)\) per path (plus the \(O(\log n)\) for the LCA).
Memory: \(O(n\log n)\) for the second table.
\end{complexitybox}

% ---------------------------------------------------------------
\subsection{First-Ancestor-Past-a-Threshold (Bit Descent)}
% ---------------------------------------------------------------

A second pattern falls out of the same table.  Suppose edges have
weights and you want the \emph{deepest ancestor} of~\(v\) whose
prefix-sum-from-root is still \(\le T\).  Equivalently: the highest~\(u\)
such that \(\dist(u, v) \le \text{cap}\) with some monotone predicate.
Use the \(\code{up}\) table in \emph{descending} bit order:
\[
\text{for } j = \mathit{LOG}\!-\!1 \text{ down to } 0:\quad
\text{if } \code{up}[\mathit{cur}][j] \text{ exists and } P(\code{up}[\mathit{cur}][j]) \text{ is false}:
\quad \mathit{cur} \leftarrow \code{up}[\mathit{cur}][j].
\]
After the loop, \(\code{par}[\mathit{cur}]\) is the first ancestor passing~\(P\).
This mirrors the lifting-LCA ``descend then return parent'' pattern.

\begin{insightbox}[One loop shape, three uses]
Descending bits + conditional jump is the \emph{same} structure as
the LCA descend loop and the binary-search-on-a-tree technique.  Once you
write it for LCA, the threshold variant is a predicate swap.
\end{insightbox}

% ---------------------------------------------------------------
\subsection{The k-th Node on a Route}
% ---------------------------------------------------------------

Given \(u, v\) and \(k\) (\(1 \le k \le \dist(u,v)\)), find the
\(k\)-th node on the path from~\(u\) to~\(v\):
\[
  d_1 = \depth(u) - \depth(\ell), \qquad d_2 = \depth(v) - \depth(\ell).
\]
If \(k \le d_1\): the answer is the \(k\)-th ancestor of~\(u\)
(just \(\texttt{kth\_ancestor}(u,k)\)).
If \(k > d_1\): the answer is on the \(v\to\ell\) arm at position
\(d_1 + d_2 - k + 1\) from~\(v\) downward from~\(\ell\), i.e.\
\(\texttt{kth\_ancestor}(v,\; d_1+d_2-k)\).

\begin{tracebox}[Route positions on T: u=8, v=10]
Path: \(8 \to 5 \to 2 \to 1 \to 3 \to 7 \to 9 \to 10\), length 7 edges.
\(\ell = 1\), \(d_1 = 3\), \(d_2 = 4\).

\(k=1\): ancestor-1 of 8 = 5.  \checkmark
\(k=3\): ancestor-3 of 8 = 1.  \checkmark
\(k=4\): \(d_1+ d_2 - k = 3+4-4 = 3\); ancestor-3 of 10 = 7?
\(10\to9\to7\to3\): 3rd ancestor = 3.  Hmm.
\(k=4\) is 1 edge past \(\ell\): the path position is \(d_1(=3) + 1 = 4\)
means the \((k-d_1)\)-th descendant of~\(\ell\) on the \(v\)-arm.
Equivalently: \(d_1+d_2 - k = 3\) ancestor steps from~\(v=10\):
\(10 \xrightarrow{1} 9 \xrightarrow{2} 7 \xrightarrow{3} 3\).  Node~3.  \checkmark

\(k=5\): \(d_1+d_2-k = 2\); ancestor-2 of 10 = 7.  \checkmark
\(k=7\): \(d_1+d_2-k = 0\); ancestor-0 of 10 = 10 itself.  \checkmark
\end{tracebox}

% ---------------------------------------------------------------
\subsection{Problem Lab}
% ---------------------------------------------------------------

\solvedproblem{c11-1304e}{1-Trees and Queries}{CF 1304E}
  {https://codeforces.com/problemset/problem/1304/E}{CF 2000 L5}
  {binary lifting, distances, parity, k-th node on route}

\begin{problemstatement}[CF 1304E 1-Trees and Queries]
A tree of \(n\) vertices, plus \(q\) queries.  Each query gives \(a, b, x,
y, k\): you must walk from~\(a\) to~\(b\), traversing exactly \(k\) edges.
You \emph{may} add edge \((x,y)\) (temporarily, just for this query) and
use it \emph{at most once}.  Can you arrive at~\(b\) in exactly \(k\) steps?
(You may revisit vertices/edges; only parity and minimum length matter.)

\textbf{Input.} \(n\) then \(n\!-\!1\) edges then \(q\) then queries.
\(2\le n,q\le 10^5\).

\textbf{Example.}
\begin{verbatim}
Input:
5
1 2
2 3
3 4
4 5
5
1 3 1 2 2
1 4 1 3 2
1 4 1 3 3
4 2 3 3 9
5 2 3 3 9

Output:
YES
YES
NO
YES
NO
\end{verbatim}
\end{problemstatement}

\begin{attemptgate}
\peek{Without the extra edge, the shortest \(a\to b\) route has
\(\dist(a,b)\) edges.  Any walk of length \(k\) must satisfy
\(\dist(a,b) \le k\) and \((k - \dist(a,b))\) even.}
\nudge{The extra edge \((x,y)\) gives two alternative route lengths:
\(\dist(a,x)+1+\dist(y,b)\) and \(\dist(a,y)+1+\dist(x,b)\).  Check all
three (direct + two detours) against the parity condition.}
\fullhint{YES iff any of the three lengths \(d_1 = \dist(a,b)\),
\(d_2 = \dist(a,x)+1+\dist(y,b)\), \(d_3 = \dist(a,y)+1+\dist(x,b)\) satisfies
\(d_i \le k\) and \((k - d_i) \% 2 == 0\).}
\end{attemptgate}

\solstep{Understanding}
The walk is on a tree (one route between every pair).  Adding edge \((x,y)\)
creates exactly one cycle; using it \emph{at most once} means the walk
goes through \((x,y)\) in one direction.  The rest is free movement (back
and forth on any edge), so you can always waste an \emph{even} number of
extra steps.  The condition is: minimum route length \(\le k\) and parity
match.

\solstep{Key Observations}
\begin{observation}
The edge must be used \emph{exactly once}, so only two candidate walk
lengths matter:
\(d_2 = \dist(a,x)+1+\dist(y,b)\) and
\(d_3 = \dist(a,y)+1+\dist(x,b)\).
Each is computed via the distance-through-hinge formula.
\end{observation}
\begin{observation}
Back-and-forth on any edge wastes exactly 2 steps, so \((k - d_i)\) must
be even and \(\ge 0\).  The answer is YES if \emph{any}~\(d_i\) works.
\end{observation}

\solstep{Think It Through}
\begin{thinkbox}
\deadend{BFS per query over the augmented graph}
{\(O(n)\) per query; \(O(nq)\) total.  TLE.}
\aha{Precompute LCA + distances in \(O(n\log n)\).  Each query is three
distance computations (each \(O(\log n)\)) + a parity check.}
\end{thinkbox}

\solstep{Implementation}
\begin{lstlisting}[caption={CF 1304E full solution (lifting distances + parity)}, label={lst:c11-1304e-soln}]
#include <bits/stdc++.h>
using namespace std;
const int LOG = 17, MX = 100001;
vector<int> adj[MX];
int up[MX][LOG], dep[MX], tin[MX], tout[MX], tmr;

void dfs(int v, int p, int d) {
    tin[v] = ++tmr; dep[v] = d; up[v][0] = p;
    for (int j = 1; j < LOG; j++) up[v][j] = up[up[v][j-1]][j-1];
    for (int c : adj[v]) if (c != p) dfs(c, v, d+1);
    tout[v] = tmr;
}
bool is_anc(int u, int v){ return tin[u]<=tin[v]&&tin[v]<=tout[u]; }
int lca(int u, int v) {
    if (is_anc(u,v)) return u; if (is_anc(v,u)) return v;
    for (int j=LOG-1;j>=0;j--) if(!is_anc(up[u][j],v)) u=up[u][j];
    return up[u][0];
}
int dist(int u, int v){ int l=lca(u,v); return dep[u]+dep[v]-2*dep[l]; }

int main(){
    int n; scanf("%d", &n);
    for (int i=0;i<n-1;i++){int a,b;scanf("%d%d",&a,&b);adj[a].push_back(b);adj[b].push_back(a);}
    dfs(1,1,0);
    int q; scanf("%d", &q);
    while (q--) {
        int a,b,x,y,k; scanf("%d%d%d%d%d",&a,&b,&x,&y,&k);
        int d2 = dist(a,x)+1+dist(y,b);
        int d3 = dist(a,y)+1+dist(x,b);
        bool ok = false;
        if (d2 <= k && (k-d2)%2==0) ok = true;
        if (d3 <= k && (k-d3)%2==0) ok = true;
        puts(ok ? "YES" : "NO");
    }
}
\end{lstlisting}

\solstep{Trace on the Official Sample}
\begin{tracebox}[CF 1304E: chain 1--2--3--4--5; all 5 queries verified]
The edge \((x,y)\) MUST be used exactly once.  Two candidate walk lengths:
\(d_2 = \dist(a,x)+1+\dist(y,b)\) and \(d_3 = \dist(a,y)+1+\dist(x,b)\).
YES iff \(d_i \le k\) and \((k - d_i)\) is even.

Depths: \(d = (0,1,2,3,4)\).  On a chain, \(\text{lca}(a,b) = \min(a,b)\),
so \(\dist(a,b) = |a-b|\).

Q1\;(1\;3\;1\;2\;2): \(d_2=|1{-}1|+1+|2{-}3|=0{+}1{+}1=2\);
\(d_3=|1{-}2|+1+|1{-}3|=1{+}1{+}2=4\).
\(k{=}2\): \(d_2{=}2\le 2\), \((2{-}2){=}0\) even \(\Rightarrow\) YES. \checkmark

Q2\;(1\;4\;1\;3\;2): \(d_2=|1{-}1|{+}1{+}|3{-}4|=0{+}1{+}1=2\);
\(d_3=|1{-}3|{+}1{+}|1{-}4|=2{+}1{+}3=6\).
\(k{=}2\): \(d_2{=}2\le 2\), \((2{-}2){=}0\) even \(\Rightarrow\) YES. \checkmark

Q3\;(1\;4\;1\;3\;3): \(d_2=2\); \(d_3=6\).
\(k{=}3\): \(d_2{=}2\le 3\), \((3{-}2){=}1\) odd NO; \(d_3{=}6>3\) NO. \checkmark

Q4\;(4\;2\;3\;3\;9): \(d_2=|4{-}3|{+}1{+}|3{-}2|=1{+}1{+}1=3\);
\(d_3\) same (\(x{=}y{=}3\)).
\(k{=}9\): \(d_2{=}3\le 9\), \((9{-}3){=}6\) even \(\Rightarrow\) YES. \checkmark

Q5\;(5\;2\;3\;3\;9): \(d_2=|5{-}3|{+}1{+}|3{-}2|=2{+}1{+}1=4\);
\(d_3=4\) (same, \(x{=}y\)).
\(k{=}9\): \(d_2{=}4\le 9\), \((9{-}4){=}5\) odd NO; \(d_3\) same NO. \checkmark
\end{tracebox}

\solstep{Complexity}
\begin{complexitybox}
Preprocessing: \(O(n \log n)\) (lifting table + tin/tout via DFS).
Per query: \(O(\log n)\) (two distance computations, each needing an LCA).
Total: \(O((n+q)\log n)\).  \textbf{AC.}
\end{complexitybox}

% ---------------------------------------------------------------
\subsection{Common Pitfalls}
% ---------------------------------------------------------------

\begin{warningbox}[The hinge value must be folded in once]
The \(\code{agg}\) accumulation on each arm \emph{excludes} the landing
vertex (the hinge).  If you forget \(\mathit{val}(\ell)\), the answer
misses one vertex.  If you include it on both arms, you double-count.
The template \(\texttt{lst:c11-pathagg}\) seeds \(\mathit{res} =
\mathit{val}[\ell]\) before the two loops -- never change that.
\end{warningbox}

\begin{warningbox}[Non-idempotent ops + the two-accumulator path]
The overlapping-block trick (\(\S\)\ref{sec:c11-euler}) needs
\(\oplus\) idempotent.  The hinge-split \emph{two-arm} technique does
NOT, but it also does not support \(\oplus\) that is only
\emph{commutative} (not associative).  If \(\oplus\) is non-commutative,
you must combine the two arms in the correct order: left arm accumulated
bottom-to-top, then \(\mathit{val}(\ell)\), then right arm accumulated
bottom-to-top \emph{reversed} (i.e., read top-to-bottom).  This is
relevant for string concatenation and matrix multiplication.
\end{warningbox}

\begin{warningbox}[CF 1304E: the extra edge is MANDATORY]
Many early submissions also check \(d_1 = \dist(a,b)\) as a candidate.
The problem requires traversing the new edge \emph{exactly once}; you
cannot skip it.  The only valid walk lengths involve the \(+1\) crossing.
\end{warningbox}

% ---------------------------------------------------------------
\subsection{Variations and Twists}
% ---------------------------------------------------------------

\begin{itemize}[nosep]
  \item \textbf{XOR on a path.}  \(\oplus = \mathit{xor}\); identity is
        0; not idempotent (\(x \oplus x = 0\)), so the sparse-table
        overlap trick would \emph{fail} -- but the two-arm method is fine.
        Verify: path \(4\to8\) has nodes \((4,2,5,8)\), values
        \((1,3,8,2)\); xor = \(1\oplus3\oplus8\oplus2 = 12\)?  Actually
        \(1\oplus3=2, 2\oplus8=10, 10\oplus2=8\).  XOR=8.  \hfill\lv{4}
  \item \textbf{Sum along a path.}  Same as the distance formula when
        edge weights are involved; \(\oplus = +\); non-idempotent, two-arm
        essential.  For path \(4\to8\) values \((1,3,8,2)\), sum=14.  \hfill\lv{3}
  \item \textbf{Count of vertices with property \(P\) along a path.}
        \(P\) is typically "value \(\le T\)"; requires offline + persistent
        or merge-sort tree (Ch.~21), not just a lifting table.  \hfill\lv{6}
  \item \textbf{Dynamic updates.}  When edge values change between
        queries, \(\code{agg}\) goes stale.  The upgrade is
        heavy-light + segment tree (Ch.~22) or Euler-tour + BIT (Ch.~18).
        CSES~2134 is the canonical problem.  \emph{Deferred by design.}
        \hfill\lv{6}
\end{itemize}

% ---------------------------------------------------------------
\subsection{Practice Corner}
% ---------------------------------------------------------------

\begin{practiceset}
  \practice{Path Queries II (dynamic path max)}{CSES 2134}
    {https://cses.fi/problemset/task/2134}
    {CSES L5}
    {maximum node value on a path \emph{with} point updates: the
     \(\code{agg}\) table goes stale the moment a value changes, so the
     static two-arm accumulate of this section no longer applies.  The
     intended tool is heavy-light \,+\, segment tree (Ch.~22 / Ch.~20).
     \emph{Deferred by design} -- this is the dynamic boundary the
     chapter deliberately leaves.}
\end{practiceset}

\paragraph{Back-cites.}
Ch.~13 (\(\thmref{c10-subtree-template}\)) -- the fold is the same
``combine children'' idea; here the children are replaced by the two
arms of a path.
Ch.~19/20 (\(\chref{c13}\)/\(\chref{c15}\)) -- dynamic path updates:
forward-cite only; nothing solved here survives a value change.
Ch.~22 (\(\chref{c16}\)) -- HLD decomposes a tree into chains so segment
trees handle path-update; the static subsumes to what \(\code{agg}\) does
here, but faster to code.

% ---------------------------------------------------------------
\begin{recallbox}[Section 16.4 in 10 seconds]
\(\code{agg}[v][j]\) carries an associative payload on each \(2^j\) jump.
Path queries split at the hinge (\(\lca\)), accumulate each arm via bit
decomposition, then fold in \(\mathit{val}(\ell)\).  Cost \(O(\log n)\)
per path query.  Idempotence is NOT required (unlike the sparse-table
overlap trick).  Same table powers threshold descents (``first ancestor
past \(T\)'') and \(k\)-th-node-on-route.  CF~1304E: add one edge
\((x,y)\) to a tree; a walk of exactly~\(k\) edges from~\(a\) to~\(b\)
exists iff \(\dist(a,x)+1+\dist(y,b) \le k\) (or the \(a\)-\(y\)-\(x\)-\(b\)
variant) and \((k - d)\) is even.
\selftest{
  \item Write the recurrence for \(\code{agg}[v][j]\).
  \item Why must you merge \(\mathit{val}(\ell)\) separately?
  \item Max along path \(8\to10\) on~\(T\) with the given values?  (9.)
  \item CF~1304E: why is \(d_1 = \dist(a,b)\) not a valid candidate?
  \item Give the \(k\)-th-node-on-route formula for \(k > d_1\).
}
\end{recallbox}
